% Track N3 — candidate section for the report.
% Numbers from results/n3/{n3a,n3b,n3c,n3d}/summary.json at the commits in docs/N3_STATE.md.
% Pre-registrations: configs/n3a_anchor_prereg.yaml (+ amendments A and B),
% configs/n3b_anchored_tiers_prereg.yaml, configs/n3c_drift_flag_prereg.yaml.
% Do not edit the main LaTeX files from this branch.

\section{Can the start attitude come from the path instead of the pixels?}
\label{sec:n3-anchored-start}

\paragraph{Decision.}
\textbf{No: the anchor is the most accurate start in the study and it is almost never available.}
A quiet window in the trajectory pins the start by physics, and where one exists the anchored lean
is accurate to a median $0.03^\circ$ on feasible clips --- within $0.01^\circ$ of an oracle start and
about a degree better than the reader --- but the window exists in \textbf{9 of 200} simulated
development clips ($4.5\,\%$), in \textbf{0 of 131} dynamic clips, and in \textbf{0 of 14} UZH
sequences and \textbf{0 of 8} TII flights once rests on a support are excluded. The pre-registered
availability gate fails, the anchored arm detects 6/27 tilt removals at a matched 8/80 pooled
false-reject against the reader's 22/27, and no forest removal and none of the six drift-tail clips
has an anchor at any threshold up to $6^\circ$.

\paragraph{Question.}
Track N and N2 established that attitude propagates accurately from video rotation given a start
good to about $2^\circ$, and that the reader's start imports a scene bias of $8$--$10^\circ$ in the
simulated forest. The start was therefore the whole problem. This section tests a start that needs
no pixels: when acceleration and velocity are near zero, differential flatness forces a zero lean,
so the thrust axis is vertical and the world-up in the camera is the known mount rotation applied to
vertical, $u_{\mathrm{cam}}(t_a) = R_{cb}^{\top}e_3$; the existing propagation then runs forward and
backward from that frame.

\paragraph{Setup.}
The anchor statistic is the implied tilt of the thrust axis,
$\theta_{\mathrm{imp}}(t) = \angle(a + g e_3 + k_d v,\; e_3)$, with the frozen values
$g = 9.80665$ and $k_d = 0.295\,\mathrm{s}^{-1}$; on the simulated truth it reproduces the stored
required tilt to $0.005^\circ$. Three tiers were pre-registered from the physics ---
$\theta_{\mathrm{imp}} \le 2^\circ$ held for $1\,\mathrm{s}$ (primary), $\le 4^\circ$, $\le 6^\circ$ ---
together with a selection rule (smallest mean implied tilt, ties earliest) and a definition of the
anchor error. Part A uses truth trajectories only: the 200 development clips at $250\,\mathrm{Hz}$,
the 14 UZH-FPV sequences at $500\,\mathrm{Hz}$ and the eight TII development flights at
$274.73\,\mathrm{Hz}$. Part B propagates DPVO rotation from the anchor and runs $F_0$, $C_2$ and the
coherence tier against the reader at a matched pooled false-reject with abstentions counted as
non-detections. Part C pre-registers a label-free reliability flag for the $20^\circ$ drift tail.
No held-out clip was read; compute was CPU only and cost nothing.

\paragraph{Anchor availability and accuracy (Part A).}
Table~\ref{tab:n3-availability} is the central result. A $2^\circ$, $1\,\mathrm{s}$ anchor exists in
the eight hover clips plus one accel-brake clip and in nothing else: every manoeuvring family
(ellipse, turn, slalom, lemniscate, climb-turn, accel-brake) is at zero. The full availability
surface shows the trade-off is not gentle: reaching a third of the clips needs a $10^\circ$
allowance at $1\,\mathrm{s}$, and a $10^\circ$ start error is the size of the forest bias the reader
already has. The quietest $1\,\mathrm{s}$ window of a development clip has a maximum implied tilt of
$3.6^\circ$ at the 10th percentile and $15.4^\circ$ at the median.

\begin{table}[t]
\centering
\caption{Anchor availability ($k/n$) and anchor error by tier. ``Dynamic'' is the Track I
definition, peak required tilt $\ge 15^\circ$ (131 of 200 clips). The anchor error is the angle
between the true world-up in the camera at the anchor frame and $R_{cb}^{\top}e_3$, which equals the
clip's true tilt there and contains no mount term.}
\label{tab:n3-availability}
\begin{tabular}{lrrrrr}
\hline
tier & all 200 & dynamic 131 & forest feas.\ 26 & err.\ median & err.\ p90 \\
\hline
$2^\circ$, $1\,\mathrm{s}$ & 9/200 & \textbf{0/131} & 3/26 & $0.00^\circ$ & $0.00^\circ$ \\
$4^\circ$, $1\,\mathrm{s}$ & 25/200 & 0/131 & 4/26 & $2.24^\circ$ & $2.81^\circ$ \\
$6^\circ$, $1\,\mathrm{s}$ & 37/200 & 2/131 & 4/26 & $2.30^\circ$ & $3.54^\circ$ \\
\hline
UZH, 14 sequences, $2^\circ$ & 0/14 & 0/14 & --- & --- & --- \\
TII, 8 flights, $2^\circ$ & 0/8 & 0/8 & --- & --- & --- \\
\hline
\end{tabular}
\end{table}

On real footage the physics needs one correction that the simulated set does not: while the drone
rests on a support a normal force replaces thrust, so a quiet window there does not pin the start.
All twelve raw UZH anchors sit before the first sustained motion, with a true tilt of
$0.33$--$9.24^\circ$ while the implied tilt is at most $0.12^\circ$, and the raw TII anchors are
pre-start rests or windows after the last sustained motion. Requiring the window to lie between the
first and the last $|v| \ge 2\,\mathrm{m\,s^{-1}}$ removes them all.

\paragraph{Verdict tiers from the anchor (Part B, descriptive because the availability gate fails).}
At the matched pooled $8/80$ feasible false-reject the anchored arm flags 8 of the 16 feasible clips
it can score and detects 6/27 tilt removals (reader 22/27), 6/27 mirrored rolls (15/27) and 0/9
forest removals (9/9); the six drift-tail clips and all nine forest removal clips have no anchor at
any tier, so those cells are fixed by coverage, not by the statistic. The informative reading is on
the clips that do have an anchor: there the anchored $F_0$ is a median $0.03^\circ$ on feasible
clips against the reader's $1.17^\circ$ (maximum $9.08^\circ$) and a reader start's $1.08^\circ$, and
a cut at the worst anchored feasible residual separates $0/16$ feasible from 6/6 removals, 5/6
mirrors and 6/6 scale edits. Those cells hold 16 and 21 clips and their null is a simulated hover
with exactly zero lean, so they are a mechanism check rather than a performance claim.

\paragraph{The drift tail is half predictable (Part C).}
The tail $F_0$ cut is set by six feasible clips above $20^\circ$, all at $30$ or $45^\circ$ mount. A
single-draw rotation-increment roughness (the 90th percentile of $|w_k - w_{k-1}|$ on the per-frame
rotation increments) flags three of the six at a pre-registered $10\,\%$ abstention budget; abstaining
on those twenty clips lets the oracle-start cut fall from $16.44^\circ$ to $10.92^\circ$ and raises its
removals from 13/27 to 15/24 at an unchanged false-reject count (8 of 74). The pre-registered recall
criterion ($\ge 2/3$) fails, and the two-draw disagreement is the better flag ($8/8$ forest removals
after abstaining) at twice the pose cost.

\paragraph{What this establishes.}
The obstacle to a pixel-free start is coverage, not accuracy: the simulated manoeuvres and the real
flights analysed here never stop in the air, so the trajectory almost never pins the vertical where
the reader is biased. A generator or a dataset that includes genuine hovers --- delivery, inspection,
a policy that pauses --- would give the anchored lean the accuracy measured here; it would not give
it reach over the dynamic clips that the verifier exists to check. The practical reading is that the
start remains a sensor or conditioning-input question (a $2^\circ$ external attitude, as N2A priced),
not a trajectory-processing one.
